\documentclass[UTF8,11pt]{ctexart}
\usepackage[a4paper,margin=24mm]{geometry}
\usepackage{amsmath,amssymb,booktabs,graphicx}
\usepackage[hidelinks]{hyperref}
\usepackage{listings,tikz}
\usetikzlibrary{arrows.meta,positioning}
\emergencystretch=3em
\lstset{basicstyle=\ttfamily\small,breaklines=true,columns=fullflexible}
\title{确定性智能体评测中的失败归因\\
\large 比较路径与匹配合同的重放研究}
\author{匿名作者\\匿名单位}
\date{2026年10月14日}
\begin{document}
\maketitle
\begin{abstract}
智能体评测检查通过，可能源于实现正确、比较路径对故障不敏感，或实际谓词与声称性质不一致。本文对观察者证据管线和固定版本的外部辅助评分器开展回顾性、无模型审计。在原门禁实际使用的投影、上下文与请求载荷链上注入一类作者设计的隐藏状态泄漏故障，并干预门禁的输入转发。12个合成时间线族包含288个嵌套条件，其中216个含状态断言的比较臂在三个产物层均拒绝该故障。仅改变故障函数的关键字参数名称，会改变基于签名的充分性标签，却不改变失败记录。相反，在门禁中丢弃扰动输入，会使同一泄漏实现得到通过判决，九个可见信息正控仍全部通过；恢复转发后，失败恢复。外部匹配器的六个词法探针对中有三个改变命中判决，但不据此认定语义等价或上游错误。目标字段对齐和评分聚合控制进一步区分匹配行为、标签一致性与检查点后果。本文贡献是包含健康、故障和恢复控制的可复现失败归因案例，不是通用检出率估计、新的非干扰理论或语言模型质量评测。
\end{abstract}
\paragraph{关键词} 智能体评测；故障注入；非干扰测试；比较有效性；时间证据；可复现性

\section{引言}
智能体在动态环境中能够知道的内容，比世界的全局状态更窄。观察者离开房间后，房间可能已经安静；记得曾经听到声音，不等于能够断言现在仍有声音。评测既要区分观察者可获得的事实和评测者独享的事实，也要区分被测智能体的失败和测量工具的失败。用过期目标字段计算出的分数不能回答新目标的问题；没有收到干预输入的比较，也不能检验该干预的影响。

后一个区分是本研究的重点。观察者证据项目的早期解释认为，安全投影器有意忽略隐藏输入，因此隐藏孪生输出相等的检查无法检出故障。该解释混淆了健康实现和比较门禁：门禁只要转发隐藏扰动，即使安全实现忽略它，也可能拒绝读取它的替换实现。本文在实际执行链上检验这种区别，不从函数签名推断检出能力。同一项目曾将外部匹配器的词法差异描述为违反已认证表示等价性；输入材料并未建立这种认证。本文保留实测行为，收窄其解释。

研究对象是生成观察者可见投影、简化模型上下文和序列化请求载荷的确定性链。我们分别干预投影器及其调用门禁。固定版本的外部辅助评分器提供另一种归因问题，涉及词法匹配、目标标签及泄漏指标聚合。两部分均没有新增模型调用。

研究问题如下。\textbf{RQ1} 实际比较能否区分健康的隐藏状态不变性与依赖隐藏状态的投影器？\textbf{RQ2} 参数命名启发式及可见信息正控通过，是否足以证明隐藏故障比较有效？\textbf{RQ3} 没有独立语义关系或检查点后果证据时，匹配器探针对能够说明什么？

贡献包括：同链泄漏注入及健康、签名改名、输入丢弃和恢复控制；区分实现行为、干预传递和报告标签的归因分析；公开源代码提取件、全部12族的不透明衍生夹具、逐比较报告和固定评分器的合成探针。证据支持具体反例，不支持故障发生率、全面变异充分性、真实感知质量或模型排名。

\section{背景与相关研究}
\subsection{时间与观察者证据}
事实有效时间与系统获知时间的区分源于时间数据库研究\cite{snodgrass}。W3C语义传感器网络本体区分现象时间与结果时间\cite{ssn}。本文将既有概念用于观察者时间线，不将时间有效性或来源追溯作为新概念。

LongMemEval考查长期对话记忆中的时间推理、知识更新和弃答\cite{longmemeval}；TimeChara考查特定时间点的角色知识\cite{timechara}。这些任务说明可获得证据与全局信息的区分具有应用意义。本文则检验模型能力评估前的确定性产物构建。三值可回答性同样不是新概念，FEVER包含信息不足类别\cite{fever}。本实现的投影器输出current和history，unknown可回答性仍在评测侧；本文没有实现完整三路回答机制。

\subsection{非干扰、空洞性与基于故障的测试}
非干扰限制不同信息域之间的影响\cite{goguen}；自组合通过成对执行分析关系性质\cite{barthe}，非干扰测试也有直接研究\cite{hritcu}。隐藏孪生输出比较属于既有关系测试。有限执行能够给出反例，不能证明所有输入上的非干扰。

空洞性检测关注通过却未有效检验预期义务的规范\cite{beer}。模型检查中的形式定义不能无差别套用于一切弱检查。本文仅将指定输入域内恒真的谓词称为逻辑恒真；将干预在比较计算之前被丢弃称为比较输入路径失败。覆盖不足、阈值宽松和缺乏语义依据是不同问题。

变异测试考查测试对人为故障的响应\cite{jia}。本文采用这一思路，但不是代表性变异测试活动。只有一类泄漏投影语义和一类门禁故障；第二种签名包装相同泄漏语义。不同时间线比较臂是在不同输入上反复检验同一故障，不能计为独立变异体。

\subsection{蜕变关系与评分合同}
蜕变测试在缺乏传统预言机时检验相关输入输出的关系\cite{segura,chen}。关系本身仍需依据。全角数字在某些任务中可能保留人类阅读含义；插入逗号或改变周边文本可能改变含义。整词匹配也可能是有意设计。判决变化不能单独决定哪一侧正确。

因此六个作者编写的成对输入被称为词法表示探针。本文另设明确受保护值的目标对齐控制，以及明确消费模式集合的聚合控制，避免将匹配器行为上升为公开基准分数错误。本文不声称提出新的变异框架，也不声称穷尽评测工具研究。

\section{研究设计与对象}
\subsection{范围与选取}
主要对象是一个项目既有的证据评测链、封存创作语料、孪生门禁及历史报告。研究是在报告解释受到质疑后选取的；故障、探针和分析在阅读实现后设计。本文是回顾性探索研究，没有预注册、随机选取被审计对象或留出故障集。“封存”仅表示语料此前已冻结，不使后续审计成为确认性研究。

12个时间线族共享一个创作程序。九族具有状态断言及隐藏事件边界，三族为不适用该干预的通信族。每族包含三种刷新策略、两种投影策略和四个查询：
\[
12\times3\times2\times4=288\text{个嵌套条件}.
\]
隐藏状态干预适用于九个状态族：
\[
9\times3\times2\times4=216\text{个比较臂}.
\]
其余72个条件不进入该分母，既不计为检出成功，也不计为失败。时间线族是作者设计的覆盖集合，不是环境独立样本。本文只报告描述性计数，不报告显著性检验或置信区间。比较臂和产物层不能作为总体推断的独立重复\cite{hurlbert}。

\subsection{证据合同与投影}
材料区分来源登记、观察者证据流、刷新轨迹、观察者可见投影及仅评测者可用的Oracle。证据记录不透明主体、谓词、断言类型、观察时间、可选事件时间、模态、范围、参与者与支持规则。来源声明的有效区间不同于刷新策略推定的有效期。这些是结构化元数据，本研究不执行图像、音频或跨模态特征提取。

轨迹规定每个查询可读取的证据快照。刷新策略为按查询、固定间隔及观察者可见事件触发；投影策略为证据范围规则及时间戳TTL。状态依据支持规则、策略和同一槽位内的观察新旧，进入current或history。槽位由主体、谓词、目标空间和断言类型组成。观察时刻为\(t_o\)，声明截止为\(t_v\)，查询为\(t_q\)时，检查
\[
t_o\leq t_q<t_v,
\]
并要求该状态为槽位内最新观察。TTL用\(t_o+\delta\)代替\(t_v\)；没有声明区间的证据范围点观察要求\(t_q=t_o\)。计划与已见证发生事项分开。规则的存在不证明真实来源内容为真。

投影器排除未发给观察者的通信以及不是观察者执行的自行动收据。本文不证明授权检查完备。私有空间验证器不在复现包中；公开验证器只检查重放模式、标识符、快照和摘要，是实验输入检查而非安全认证。

\subsection{实际比较链}
以\(E,R,T,P,q,H\)表示证据、来源、轨迹、策略、查询和隐藏平面：
\[
V=F(E,R,T,P,q;H),\quad C=G(V,t_q),\quad
B=J(C,\mathrm{question},\mathrm{model},\mathrm{budget}).
\]
\(V\)为投影，\(C\)为简化上下文，\(B\)为载荷。\(G\)仅暴露许可声明并去除投影元数据。\(J\)使用显式合成模型名、零温度和固定token上限构建消息，不发送请求。公开适配去除配置加载并要求显式模型参数；提取记录提供原函数及摘要。

孪生保持\(E,R,T,P,q\)不变，仅改变\(H\)。每个状态的隐藏结束时刻分别位于查询前后八分钟，起点和槽位不变。安全\(F\)有意忽略\(H\)；泄漏替身\(F_\ell\)依据隐藏区间重新分配已准入证据的current/history，再沿同一\(G,J\)继续。差异来自该依赖，而不是另一套参考上下文构建器。

\begin{figure}[t]
\centering
\begin{tikzpicture}[node distance=7mm,box/.style={draw,rounded corners,align=center,font=\small,text width=25mm,minimum height=12mm},>=Latex]
\node[box] (h) {隐藏孪生\\查询前／后};
\node[box,right=of h] (a) {门禁\\转发／丢弃};
\node[box,right=of a] (f) {投影器\\安全／泄漏};
\node[box,right=of f] (g) {上下文\\请求载荷};
\draw[->] (h)--(a);\draw[->] (a)--(f);\draw[->] (f)--(g);
\node[font=\small,align=center,below=8mm of f] (e) {固定证据、来源、快照、查询和策略};
\draw[->] (e)--(f);
\end{tikzpicture}
\caption{实际比较路径。安全排除隐藏输入是预期行为；同链泄漏替换用于检验转发。门禁丢弃输入可以掩盖该故障，而可见信息正控仍通过。}
\end{figure}

\section{审计步骤}
\subsection{干预与控制}
表\ref{tab:design}列出六种配置。A使用健康投影器和原转发门禁；B仅替换为\(F_\ell\)；C用关键字字典名为\texttt{\_hidden\_world}的函数包装同样泄漏。历史签名启发式将该拼写用作排除隐藏输入的依据，但名字不是数据流证明：C会读取并转发这些值。

D保留泄漏实现，在门禁加入接受却丢弃隐藏参数的适配器；E对安全投影器使用同一适配器。E中的输出相等符合预期，但不能证明门禁健康。F保留泄漏投影器并移除适配器。F撤销的是作者注入故障，不是修复原转发门禁中自然观察到的故障。

历史判决逻辑保持不变。\texttt{PASS\_REFERENCE\_ONLY}的退出码为1，不能单看退出码判断实现失败。外层审计可正常完成而嵌套门禁报告FAIL；报告区分两种状态。

\begin{table}[t]
\centering\small
\caption{声明的干预。B/C为同一泄漏语义；E/F是补充控制。}
\label{tab:design}
\begin{tabular}{@{}clll@{}}
\toprule
编号 & 投影器 & 转发路径 & 作用\\
\midrule
A & 原安全实现 & 原门禁 & 健康实现\\
B & 泄漏替身 & 原门禁 & 实现故障\\
C & 相同泄漏、参数改名 & 原门禁 & 签名标签控制\\
D & 泄漏替身 & 丢弃适配器 & 比较路径故障\\
E & 原安全实现 & 丢弃适配器 & 健康输出、故障路径\\
F & 泄漏替身 & 恢复原转发 & 撤销人工路径故障\\
\bottomrule
\end{tabular}
\end{table}

\subsection{测量与停止规则}
每个适用臂比较投影、上下文和载荷，分别记录不相等臂数、族数及层级条目。同臂三个层级失败是同一故障的关联后果。九个可见信息正控改变观察者证据，经过不同输入路径；其结果和签名标签保留。

公开重放新增测量元数据，不替换历史判决。对人工丢弃路径的无效标记来自已知干预分配，不是自动发现未知故障。适用性、实际相等性、注入故障和历史状态保持独立字段。

全部六配置和声明比较臂执行后停止；输入不合法则失败。不会继续寻找故障直到出现期望结论。冻结输出目录禁止覆盖，输入和历史记录保持原状。

\subsection{外部评分器控制}
GateMem辅助评分器固定于提交\texttt{603f9f4b4ba4b77f043c20f85687fa016fd720b0}。其说明明确LLM评审是主评估、规则评分为辅助。本文不评测其LLM评审或整个基准。MIT许可的\texttt{metrics.py}、\texttt{scorer.py}逐字节保留，并附许可和摘要。

六个作者设计的输入对直接调用\texttt{pattern\_hits}。大小写和空白探针保持数字目标；全角和逗号探针同步改变目标；CJK探针改变周边文本但保留ASCII目标。逗号例不是规范千位分组，CJK文本也不被认定为等义翻译。不认证任何语义变换类别。

隐私控制明确合成受保护PIN，令上下文由2864变为9999。过期标签分支在\texttt{leak\_targets}及\texttt{judge\_spec.not\_include}中保留2864；对齐分支更新两者，消费相同的改变后预测字节。最终回答拒绝披露，不含两种PIN。上游\texttt{score\_one}分别报告回答、上下文和端到端泄漏。

聚合控制的上下文为\emph{reference 99990 recorded marker}。比较仅含9999的模式集合与额外加入marker的集合。该控制检验模式析取，不判定数字前缀是否应命中；未使用上游真实episode。

\section{结果}
\subsection{RQ1：原转发比较能够拒绝注入泄漏}
A的隐藏产物相等，历史状态为\texttt{PASS\_REFERENCE\_ONLY}。B在九族中出现216个不相等臂；每臂投影、上下文和载荷都不同，合计648个层级条目。两配置没有其他报告失败，九个可见信息正控均通过。

该结果修正了仅因安全投影器忽略输入，就认定门禁无法检出隐藏依赖的推断。健康相等符合合同。B证明原转发门禁能沿被测链拒绝这一替换故障，不证明A的全面非干扰、其他故障检出或通用变异充分性。216次机会检验的是一个故障，不是216个变异体。

\begin{table}[t]
\centering\small
\caption{同链结果。每行288条件、216适用比较。RF为历史REFERENCE\_ONLY签名标签，EC为EVALUATED\_CHAIN；均不是证明。}
\begin{tabular}{@{}clrrrrl@{}}
\toprule
编号 & 历史状态 & 不相等臂 & 层级条目 & 受影响族 & 正控 & 标签\\
\midrule
A & PASS\_REFERENCE\_ONLY & 0 & 0 & 0 & 9 & RF\\
B & FAIL & 216 & 648 & 9 & 9 & EC\\
C & FAIL & 216 & 648 & 9 & 9 & RF\\
D & PASS & 0 & 0 & 0 & 9 & EC\\
E & PASS\_REFERENCE\_ONLY & 0 & 0 & 0 & 9 & RF\\
F & FAIL & 216 & 648 & 9 & 9 & EC\\
\bottomrule
\end{tabular}
\end{table}

\subsection{RQ2：签名和可见正控不是充分条件}
C与B的不相等臂和层级条目相同，却获得REFERENCE\_ONLY。参数命名改变标签，没有改变泄漏语义或实际拒绝。因此该启发式不能可靠描述这些实现是否使用隐藏输入。

D返回PASS、零不相等臂及EVALUATED\_CHAIN；泄漏函数签名仍在，干预却未到达调用。九个可见正控全通过。该配置反驳
\[
\text{可见信息正控通过}\Longrightarrow\text{本隐藏故障比较有效}
\]
及“签名标签足以证明有效”的推断。结论针对输入路径：不同路径上的成功不能单独证明隐藏输入得到转发。

E说明只看输出相等也无法诊断门禁；健康实现在已知丢弃路径上仍相等。F移除适配器后恢复216个失败臂。干预与恢复明确了本设置中D掩盖故障的原因，不证明其自然发生率。

\subsection{RQ3：词法行为不是语义错误}
六对中三对改变命中判决。实现根据归一化目标是否属于某ASCII型字符类选择分支：属于者使用Unicode单词边界，否则用子串匹配。全角数字和逗号改变分支；CJK邻接改变ASCII目标的单词边界；大小写和空白例在选定输入上保持判决。

\begin{table}[t]
\centering\small
\caption{词法探针结果。F/T为未命中/命中。JSON含精确Unicode字串，均无语义等价认证。}
\begin{tabular}{@{}p{.19\linewidth}p{.52\linewidth}cc@{}}
\toprule
探针 & 基线到变体 & 基线 & 变体\\
\midrule
全角 & 99990/9999变为全角文本及目标 & F & T\\
逗号 & 99990/9999变为9,9990/9,999 & F & T\\
CJK邻接 & 空格分隔9999变为CJK紧邻9999，目标不变 & T & F\\
上下文大写 & reference变为REFERENCE，目标不变 & T & T\\
ASCII空格 & 单空格变重复空格，目标不变 & T & T\\
全角空格 & ASCII空格变U+3000，目标不变 & T & T\\
\bottomrule
\end{tabular}
\end{table}

这些是事后选择的行为观察，不是依据共同语义合同抽样的错误率。三对变化不意味着半数Unicode变换错误，而说明匹配分支敏感，需明确评分器承诺保持的关系。

目标对齐控制中，原PIN的上下文与端到端泄漏为真；改变上下文但保留旧标签后为假；同步标签后恢复为真。回答侧始终为假。过期和对齐分支读取相同改变后预测，差异来自标签一致性，不是上游无法检测指定目标。

聚合控制中，仅含数字模式时为假，加入匹配marker后为真。因此单个字串未命中不能直接等同于检查点判决改变。控制不裁定前缀语义，明确区分匹配行为与聚合后果。

\section{解释与启示}
\subsection{采用归因记录而非不可反驳分类学}
归因记录应区分声称性质、许可输入域、目标故障、实际谓词和执行路径，并明确适用性与聚合层级。这是案例支持的报告建议，不是已证明优于其他审计流程的方法。

\(x=x\)是逻辑恒真；\(k>0\)不是，因为\(k=0\)可反驳，但它可能不足以满足覆盖要求。空集上的全称谓词为真；仅当任务还要求非空时才违反合同。健康隐藏相等是预期，门禁丢弃干预后的相等则不能说明干预影响。参数名称不能证明未读取数据，词法差异不能证明语义错误。

早期稿将这些情形归为九种不同的不可反驳机制。本文撤回该分类学及仅由健康实现推断门禁无能的主张，也撤回六探针的语义认证表述。旧稿作为历史保留，不用于支持修订后的结论。

\subsection{评测工具也是被测系统的一部分}
直接门禁干预提供了报告布尔值变异不能提供的证据。改变汇总字段可检查判决连线，但不能证明输入到达投影器。可见正控仅证明对自身变化敏感；同路径实现故障为隐藏比较能失败提供见证。

为避免循环验证，公开包展示适配器、分配、完整失败和恢复结果。有效性元数据并不自动发现未知故障。历史状态及签名字段仍保留，即使本文解释不同，读者也能检查计算而不是只接受更正标签。

\subsection{对时间敏感智能体研究的意义}
观察者与评测者平面分离使干预可定位。隐藏事件真值不能替代模型可见投影里的来源声明有效期；历史保留不能自动支持当前断言。但结构化重放不证明模型遵守这些边界。

跨模态获取、空间对齐、来源冲突、长上下文压力及语义回答评分均需新实验和声明合同。它们可以复用本底座，但不是现有计数已证明的成果。

\section{有效性威胁}
\paragraph{构念有效性。}
故障通过隐藏区间决定current/history，保留已准入集合，不插入秘密字串。它在指定合同下属于非法依赖，不是全面隐私泄漏模型。收件人、空间和来源有效性故障尚未测试。词法探针没有独立任务级等价预言机。

\paragraph{内部有效性。}
进程内确定性替换保持观察者输入和后续链不变，F明确移除人工适配器。提取及不透明化可能引入差异；记录含源文件/函数摘要、288臂current/history归属保持及四主配置报告对应检查。这证明归属和失败结构保持，不证明匿名化前后载荷逐字节相同。标识符与查询改写改变字节；衍生收据属于作者侧证明，外部读者没有原件时不能独立证明重命名。

\paragraph{外部有效性。}
仅研究一个作者设计管线和一个辅助评分器提交。共同创作及非随机选取限制推广。反例反驳被检验的充分性推断，不估计其在其他系统中的发生率。不声称首次发现或全面综述。

\paragraph{统计有效性。}
216臂和三个层级执行同一故障；648条目不是独立变异体。不估计检出概率、显著性、策略效果或模型差异。

\paragraph{历史与未完成工作。}
旧回答批次记录288调用、217可用、68空输出和三个非空截断；68空输出包含在71个长度终止中。这些是交付计数，不是语义分数，公开包不从原始回答重算。历史来源还记录后续变动的脚本摘要。准备度报告仍有84个占位确认查询。它们是背景，不支持新的确定性结果或模型表现主张；有限案例研究不要求完成这些未完成模型实验。

\section{复现与材料可用性}
公开仓库为\url{https://github.com/puresky271/observer-evidence-study}。包含双语TeX/PDF、提取函数、12匿名族、六配置完整报告、精确探针、评分源码及许可。发布清单使用实测SHA-256，提取收据标明原文件和函数切片。重放不需要网络、密钥、生产状态或模型，写入新输出位置。

公开材料是适配衍生件：保持投影与比较语义，替换身份及自由文本、重算关联摘要、归一化路径元数据，要求显式模型。身份映射、私有空间验证器、原始模型回答、上游episode和隐藏目标不公开。复现范围是新确定性实验和合成控制，不是旧稿全部历史审计。

将仓库\texttt{src}加入Python路径后执行：
\begin{lstlisting}
python -m observer_study.replay --out replay_output/run1
python -m observer_study.matcher_audit --out replay_output/matcher.json
python -m unittest discover -s tests -v
\end{lstlisting}
仓库包装命令还比较重放字段和冻结报告。成功表示复现义务匹配，不表示故障实现通过。逐臂失败和输入均可检查；双语版本呈现同一研究，英文源面向投稿。

\section{结论}
原隐藏孪生门禁能够拒绝泄漏投影器，即使其健康实现有意忽略隐藏输入。签名启发式错标相同泄漏包装；丢弃输入适配器掩盖故障，九个可见正控仍通过；移除适配器恢复216失败臂。外部控制区分词法分支、标签对齐和聚合，不认定语义等价或上游基准错误。本文建立有限反例及可复现归因链，通用故障覆盖和模型能力仍需后续证据。

\section*{声明}
作者信息按作者要求暂留匿名。新增确定性研究没有收集人类被试数据，输入为合成或不透明衍生材料，私有运行记录不公开。计算辅助用于检查、实现、修订及模拟审稿，作者对核实和投稿负责。模拟评审不是期刊编辑决定。
\paragraph{工作流来源。} 引用管理与投稿前辅助审查使用了有文档记录的科学代理程序\cite{skills}。该引用说明辅助工作流，不是实证结论的依据。
\clearpage
\small
\bibliographystyle{plain}
\bibliography{references_verified}
\end{document}
